%@ pages 336-338
%@ status translated
% Chapter I, section 6: Categories with generators and exact inductive
% limits (pp. 336-338).  The section heading stands half way down p. 336; the
% text runs on to the top of p. 338 and stops at the heading of section 7.
% Transcribed from the page images; the OCR text layer in reference/ocr/ was
% not trusted.  In the definition of a morphism of inductive systems the
% original prints "f_j o u_ji = v_ji o f_i" with "o" for the composition sign
% (p. 337); both sides go from M_i to N_j, so this is correct as printed and
% is reproduced unchanged.  Proposition 6 carries a label because p. 338
% refers back to it.

\section{Categories with Generators and Exact Inductive Limits}

\begin{proposition}
If $\Cat$ is an abelian category with generators and $X$ an object of $\Cat$,
there exists a set belonging to $\Univ$ and having the same cardinality as
the set of subobjects of $X$.
\end{proposition}

Let $(X_i)_{i\in I}$ indeed be a family of generators of $\Cat$, the set $I$
belonging to $\Univ$. We know that the sum $E$ of the sets
$\operatorname{Hom}(X_i, X)$ is an element of $\Univ$. We associate to every
subobject $Y$ of $X$ the part $EY$ of $E$ which is formed of the morphisms
whose image is contained in $Y$. If $Y'$ is a subobject of $X$ not
containing $Y$, there exists an element of $E$ not belonging to
$E(Y \cap Y')$, for $(X_i)_{i\in I}$ is a family of generators of $\Cat$.

In other words, the map $Y\mapsto EY$ is an injective map from the set of
subobjects of $X$ into the set of parts of $E$. This proves the proposition.

Let us now suppose --- and this for the rest of the paragraph --- that $\Cat$
is a category with generators and inductive limits. According to \cite{grothendieck-1}, it
comes to the same thing to say that $\Cat$ is a category with generators and
that there exists a direct sum for every family of objects indexed by a set
belonging to $\Univ$. If these conditions are realised, every increasing
filtered family $(M_n)_{n\in E}$ of subobjects of an object $M$ has a least
upper bound: according to the preceding proposition, one may indeed suppose
that $E$ belongs to $\Univ$; the canonical morphisms of $M_n$ into $M$ then
define a morphism of the direct sum $\sum_{n\in E} M_n$ into $M$. The image
of this last morphism is the least upper bound of the $M_n$.

Let $I$ be a filtered ordered set, and let $(M_i, u_{ji})$ and $(N_i, v_{ji})$
be two inductive systems of morphisms of $\Cat$ indexed by $I$
($I\in\Univ$). A morphism of the first into the second is, by definition, a
family $(f_i)_{i\in I}$ of morphisms with source $M_i$ and target $N_i$ such
that $f_j\circ u_{ji} = v_{ji}\circ f_i$ if $j > i$. Thus the category of
inductive systems indexed by $I$ is defined, and this category is abelian.
Moreover, the functor $(M_i,u_{ji})\mapsto \lim\limits_{\to} M_i$ is right
exact. If this last functor is exact for every $I\in\Univ$, we say that $\Cat$
is a \emph{category with exact inductive limits}.

\begin{proposition}\label{prop:exact-inductive-limits}
If $\Cat$ is a category with generators and inductive limits, the following
assertions are equivalent:
\begin{enumerate}[label=\alph*.]
\item $\Cat$ is a category with exact inductive limits.

\item If $(P_i)_{i\in I}$ is an increasing filtered family of subobjects of
      $P$, the canonical morphism of $\lim\limits_{\to} P_i$ into $P$ is an
      isomorphism of $\lim\limits_{\to} P_i$ onto the least upper bound
      $\sup\limits_{i\in I} P_i$.

\item If $(P_i)_{i\in I}$ is an increasing filtered family of subobjects of
      $P$, and if $Q$ is a subobject of $P$, then
      \[
        (\sup P_i) \cap Q = \sup(P_i \cap Q).
      \]
\end{enumerate}
\end{proposition}

The inductive systems of an abelian category $\Cat$ with generators and exact
inductive limits are dealt with ``in the same way'' as the inductive systems
of abelian groups. We shall make much use, in what follows, of the theory of
the semisimple objects of such a category: an object $S$ of $\Cat$ is called
\emph{simple} if it is nonzero and if it contains no subobject distinct from
$S$ or from $0$. An object is called \emph{semisimple} if it is isomorphic to
a direct sum of simple objects. The theory is in every point analogous to
the theory of semisimple modules \cite{bourbaki-algebre}. Let us recall only a few results:

Let $\mathbf{S}$ be the full subcategory of $\Cat$ whose objects are the
semisimple objects of $\Cat$; let $\mathfrak{E}$ be the set of isomorphism
classes of simple objects. If $M$ is an object of $\mathbf{S}$ and if
$\lambda$ belongs to $\mathfrak{E}$, we denote by $M_\lambda$ the isotypic
component of type $\lambda$ of $M$: it is the sum of the simple subobjects of
$M$ belonging to $\lambda$. We also denote by $\mathbf{S}_\lambda$ the full
subcategory of $\mathbf{S}$ whose objects $M$ are those such that
$M = M_\lambda$. In these conditions we have the following results:

\begin{itemize}
\item If $f\colon M\to N$ is a morphism of $\mathbf{S}$, then $f(M_\lambda)$
      is contained in $N_\lambda$.

\item The functor $M\mapsto (M_\lambda)_{\lambda\in\mathfrak{E}}$ defines
      an equivalence between $\mathbf{S}$ and the product category
      $\prod\limits_\lambda \mathbf{S}_\lambda$ (see \S~8 for the
      definition of equivalences).

\item If $S$ is a simple object of $\mathbf{S}_\lambda$, the functor
      $M\mapsto \operatorname{Hom}(S, M)$ defines an equivalence between
      $\mathbf{S}_\lambda$ and the category of vector spaces [over the field
      $\operatorname{Hom}(S,S)$] whose underlying set belongs to $\Univ$.
\end{itemize}

Let us point out, to finish, that Krull--Remak--Schmidt's theorem is true in
an abelian category with generators and exact inductive limits. More
precisely, let us say that an object $M$ is \emph{indecomposable} if it is
nonzero and if every direct factor of $M$ is equal to $0$ or to $M$. We then
prove the following theorem:

\begin{theorem}
Let $\Cat$ be an abelian category with generators and exact inductive limits.
Let $(M_i)_{i\in I}$ and $(N_j)_{j\in J}$ be two families of indecomposable
objects whose ring of endomorphisms is local ($I\in\Univ$, $J\in\Univ$). If
the direct sums
\[
  \sum_{i\in I} M_i \quad\text{and}\quad \sum_{j\in J} N_j
\]
are isomorphic, there is a bijection $h$ of $I$ onto $J$ such that $M_i$ is
isomorphic to $N_{h(i)}$.
\end{theorem}

Let us recall that a ring $A$, commutative or not, is called \emph{local} if
it possesses a nonzero unit element and if the quotient of $A$ by its
Jacobson radical is a field.

\textsc{Azumaya} has proved this theorem ``in the case of modules'' \cite{azumaya-1}. The
reader is invited to take up this proof again, making use of
Proposition~\ref{prop:exact-inductive-limits} in order to avoid the reasonings
in which \textsc{Azumaya} speaks of the elements of a module.